\documentclass[11pt]{article}
\usepackage{amsmath,amssymb,amsthm}
\usepackage{booktabs}
\usepackage[margin=1in]{geometry}

\newtheorem{theorem}{Theorem}
\newtheorem{lemma}[theorem]{Lemma}
\newtheorem{corollary}[theorem]{Corollary}

\theoremstyle{definition}
\newtheorem{definition}[theorem]{Definition}

\title{Problem 748: Upside Down Diophantine}
\author{Project Euler Solutions}
\date{}

\begin{document}
\maketitle

\section{Problem Statement}

Find primitive solutions to $\frac{1}{x^2}+\frac{1}{y^2}=\frac{13}{z^2}$ with $\gcd(x,y,z)=1$, $x \le y$. Define $S(N) = \sum(x+y+z)$ over all such with $1 \le x,y,z \le N$.

Given: $S(100)=124$, $S(10^3)=1470$, $S(10^5)=2\,340\,084$. Find the last 9 digits of $S(10^{16})$.

\section{Mathematical Foundation}

\begin{theorem}[Algebraic Reformulation]
The equation is equivalent to $z^2(x^2+y^2) = 13\,x^2 y^2$.
\end{theorem}

\begin{proof}
Multiply both sides of $\frac{1}{x^2}+\frac{1}{y^2}=\frac{13}{z^2}$ by $x^2 y^2 z^2$.
\end{proof}

\begin{theorem}[Coprime Reduction]
Let $d = \gcd(x,y)$, $x = da$, $y = db$ with $\gcd(a,b)=1$. Then:
\[
z^2 = \frac{13\,d^2 a^2 b^2}{a^2+b^2},
\]
and for $z \in \mathbb{Z}^+$ we need $(a^2+b^2) \mid 13\,d^2$.
\end{theorem}

\begin{proof}
Substituting $x = da$, $y = db$ into the reformulated equation gives $z^2 = 13d^2 a^2 b^2/(a^2+b^2)$. Since $\gcd(a,b) = 1$, we have $\gcd(a^2 b^2,\, a^2+b^2) = 1$, so $(a^2+b^2) \mid 13d^2$.
\end{proof}

\begin{theorem}[Gaussian Integer Parametrization]
The equation $a^2+b^2 = 13\,t^2$ is solved via the factorization $13 = (2+3i)(2-3i)$ in~$\mathbb{Z}[i]$. Solutions with $\gcd(a,b)=1$ are parametrized by Gaussian integer multiples of $2+3i$.
\end{theorem}

\begin{proof}
In $\mathbb{Z}[i]$, $N(a+bi) = a^2+b^2 = 13t^2$, so $a+bi$ is divisible by $2+3i$ or $2-3i$. Writing $a+bi = (2+3i)\cdot w$ with $N(w) = t^2$ and applying the classification of Gaussian integers by norm yields the full parametrization.
\end{proof}

\begin{lemma}[Primitivity]
$\gcd(x,y,z) = 1$ is equivalent to $\gcd(d,z) = 1$, since $\gcd(x,y) = d$ and $\gcd(a,b) = 1$.
\end{lemma}

\begin{proof}
$\gcd(x,y,z) = \gcd(da, db, z) = \gcd(d\cdot\gcd(a,b),\, z) = \gcd(d,z)$.
\end{proof}

\begin{theorem}[Enumeration Strategy]
All primitive solutions with $x,y,z \le N$ are found by iterating over coprime pairs $(a,b)$ with $a \le b$, finding valid~$d$ such that $(a^2+b^2) \mid 13d^2$ and $z = dab\sqrt{13/(a^2+b^2)}$ is integral, then checking bounds and primitivity.
\end{theorem}

\begin{proof}
Theorem~2 reduces every primitive solution to a coprime pair $(a,b)$ and scaling factor~$d$. Systematic search over these parameters is complete by construction.
\end{proof}

\section{Algorithm}

\begin{enumerate}
\item For each coprime pair $(a,b)$ with $a \le b$ and $a^2+b^2 \le 13N^2$:
  \begin{enumerate}
  \item Compute $s = a^2+b^2$.
  \item Find all $d$ such that $s \mid 13d^2$ and $da, db, z \le N$.
  \item For each valid~$d$, verify $z^2 = 13d^2 a^2 b^2/s$ is a perfect square, check $\gcd(d,z)=1$, accumulate $x+y+z$.
  \end{enumerate}
\item Return the sum modulo $10^9$.
\end{enumerate}

\section{Complexity Analysis}

\begin{itemize}
\item \textbf{Time:} $O(N)$ with careful bounding of the $(a,b,d)$ search space.
\item \textbf{Space:} $O(\sqrt{N})$ for coprimality sieve.
\end{itemize}

\section{Answer}

\[
\boxed{276402862}
\]

\end{document}
